% !TEX root = ../thesis.tex

%=========================================================
\chapter{Conclusion}
\label{ch:conclusion}
%=========================================================

    Superconducting junctions connect macroscopic quantum coherence with microscopic charge transport. Their measured response is determined not only by the superconducting electrodes, but also by the transmission channels of the junction, the electromagnetic environment, and external microwave driving. In this thesis, I investigated how these ingredients shape quasiparticle, Cooper pair, and multiple Andreev reflection transport in aluminum nanostructures at millikelvin temperatures.

    The experiments followed a progression from oxide tunnel barriers to atomic contacts with finite and high transmission. Static transport established the superconducting energy scales, junction parameters, and, where possible, the individual channel transmissions. Microwave amplitude sweeps then revealed how the different transport contributions evolved under drive. Their voltage spacing constrained the associated charge transfer processes, while their amplitude dependence helped distinguish mechanisms with otherwise similar spectral signatures. This combination of static characterization and amplitude resolved spectroscopy provided a common strategy for testing transport models across the different junction regimes.

    This experimental strategy required a consistent measurement and analysis chain. I combined mechanically controllable break junctions with filtered low-temperature transport measurements, microwave irradiation, and synchronized acquisition of complete \textit{I--V} amplitude sweeps. I developed the digital infrastructure for instrument control, structured data storage, numerical differentiation, parameter extraction, and comparison with BCS and channel resolved MAR calculations. This common workflow connected the measured \textit{I--V}, \textit{dI--dV}, and \textit{dV--dI} responses to the transport models used throughout the thesis and made their range of validity directly testable.

    %=========================================================
    \subsubsection*{Tunnel Barriers}
    %=========================================================

        In the tunnel regime, I established a high-resolution reference for quasiparticle photon assisted tunneling. The static transport of TB1 was described by the SIS BCS model, which determined the superconducting gap and effective broadening. Pointwise fits of systematic microwave amplitude sweeps then provided an in situ calibration of the local ac voltage while the static junction parameters remained fixed. This calibration placed the complete measured transport maps on the reduced drive coordinate of the Tien-Gordon Model for direct comparison with simulation.

        The measured and simulated maps showed exceptional agreement across the complete accessible drive range. The model reproduced not only the sideband spacing, but also the redistribution of spectral weight between the main coherence peaks and the successive photon assisted replicas. Individual replicas remained visually resolved and could be followed continuously to at least order 25, forming a sideband ladder substantially wider than the superconducting gap. This simultaneous resolution of many closely spaced branches required strong microwave coupling, high voltage resolution, and a stable junction throughout the amplitude sweep. To my knowledge, no previous amplitude map of quasiparticle photon assisted tunneling in an SIS junction has resolved and continuously tracked individual replicas to a higher order under the same definition. TB1 therefore provided more than a calibration reference. It demonstrated quantitative Tien-Gordon transport deep into the multiphoton regime over an unusually wide dynamic range.

        In TB2, I observed microwave induced rectification with a local current asymmetry reaching approximately \qty{25}{\percent}. Its magnitude and polarity were tunable through three experimental controls: dc bias, microwave amplitude, and frequency. The rectification switched on above a drive threshold, responded sensitively to the microwave frequency, and reversed direction across the investigated frequency interval. This behavior turned the junction from a symmetric photon assisted tunneling reference into a microwave controlled diode-like element. Separate increasing- and decreasing-bias sweeps overlapped within the experimental resolution, which established that the response was nonhysteretic under the investigated conditions.

        Comparison with the bias symmetric Tien-Gordon Model separated the conventional photon assisted redistribution from the additional asymmetric response and established the rectification as a transport regime beyond uniform ac driving. The measurements did not determine a unique microscopic origin of this regime, and neither the local microwave field nor the voltage across the break junction was measured independently. Uniform heating could modify the quasiparticle response but could not by itself select a bias direction. The narrow frequency dependence and sign reversal instead indicated that resonant microwave coupling and a phase sensitive element were involved. I proposed that asymmetric excitation of the device leads drove the intact break junction as a Josephson weak link and thereby modified the island potential detected through the tunnel barrier. Nonequilibrium quasiparticle distributions and heating could have enhanced this response. This causal sequence connects the drive threshold, strong frequency sensitivity, and polarity reversal within a common experimentally testable picture.

        Taken together, the tunnel-barrier experiments established two complementary results. TB1 demonstrated the quantitative validity of the Tien-Gordon Model over a high-resolution sideband ladder extending deep into the multiphoton regime. TB2 revealed a tunable rectification regime beyond the model's bias symmetric assumptions. The comparison showed that nominally similar tunnel junctions can act either as controlled spectrometers of conventional photon assisted tunneling or as sensitive probes of local microwave dynamics and symmetry breaking.

    %=========================================================
    \subsubsection*{Atomic Contacts}
    %=========================================================

        Increasing the junction transmission exposed transport through discrete conduction channels. Coherent MAR calculations reproduced the static spectra along a complete opening trace and allowed me to extract the evolving mesoscopic pincode. The individual transmissions changed continuously within conductance plateaus and reorganized at abrupt conductance jumps, consistent with elastic deformation between atomic rearrangements. Channel resolved MAR spectroscopy therefore revealed changes in the contact that remained hidden in its total normal-state conductance, while constraining the contact evolution without requiring a unique reconstruction of its atomic geometry.

        Under microwave irradiation, the few-channel contacts displayed overlapping replicas of quasiparticle, MAR, and zero-bias Cooper pair transport. I developed a deterministic threshold decomposition that partitioned the simulated static MAR conductance around its known thresholds before applying charge dependent photon dressing. This construction avoided the unphysical mixing produced by direct dressing of the dc FCS components and reproduced the principal finite-bias PAMAR trajectories. Combining it with a PAICPT contribution reconstructed the complete response of C1 using the static pincode, independently determined environmental parameters, and calibrated microwave amplitude without further adjustment to the driven map. The exceptionally close agreement extended from the feature positions to the redistribution of spectral weight across the complete amplitude range.

        The amplitude dependence provided the decisive distinction between the overlapping Cooper pair mechanisms. The replicas originating from zero bias followed the squared-Bessel redistribution expected for PAICPT rather than the amplitude dependence of an ideal coherent Shapiro response. Together with the static \textit{P(E)} analysis, this identified the zero-bias contribution as incoherent Cooper pair transport. A second contact with a nearly ballistic leading channel validated the principal PAMAR description in a deliberately contrasting pincode regime. Across both contacts, threshold weighted PAMAR captured the dominant finite-bias response, while PAICPT accounted for the zero-bias contribution when it could be separated from the MAR background.

        At high conductance, the experiment entered a qualitatively different regime. For C3 at approximately \(29\,G_0\), a unique channel pincode was inaccessible and none of the tested MAR or Josephson--MAR models reproduced the sharp static characteristic convincingly. Using the measured static response as the reference nevertheless revealed an additional family of trajectories halfway between the conventional Cooper pair replicas. Their positions and squared-Bessel-like amplitude dependence consistently indicated a subharmonic PAICPT response containing an additional \(p=2\) component. The effective \(4e\) form of this component did not by itself establish correlated transfer of two Cooper pairs, but the observed subharmonic replica pattern is a new experimental result. To my knowledge, such a subharmonic PAICPT response has not been reported previously.

        Taken together, the atomic-contact measurements connected the static mesoscopic pincode to the complete microwave driven response across widely different channel compositions. They established a practical framework for separating coexisting PAMAR and PAICPT, demonstrated its reach from moderate to nearly ballistic transmission, and identified its boundary at high conductance. At that boundary, the emergence of a previously unreported subharmonic response opened a new regime for investigating charge transfer beyond the conventional few-channel description.

    Across both junction regimes, static spectroscopy provided the reference required to interpret transport under microwave drive. The voltage spacing of the driven structures constrained the effective transferred charge, while their evolution with microwave amplitude distinguished processes that produced similar spectral positions. Complete amplitude maps therefore revealed the coexistence and redistribution of transport contributions that could not be identified reliably from individual traces or feature spacing alone.

    The comparison also showed that the electromagnetic environment forms an active part of a superconducting junction experiment. It set the local microwave drive, enabled incoherent Cooper pair transfer, and contributed to symmetry breaking beyond a uniform ac excitation. The applicable transport description consequently evolved with junction transmission and environmental coupling. Importantly, the boundaries of these descriptions exposed new physics rather than only limiting the analysis: microwave induced rectification in the tunnel device and a subharmonic response in the high-conductance contact.

    %=========================================================
    \subsubsection*{Outlook}
    %=========================================================

        The microwave induced rectification provides several direct routes for further investigation. A measurement of the complex microwave transfer function at the sample would test whether the amplitude enhancement and phase evolution of a local resonance coincide with the onset and polarity reversal of the rectification. Comparing the same device before and after opening the break junction would isolate the role of the Josephson weak link. Gate and magnetic field dependent measurements within the irradiated asymmetric regime could then test whether the island potential and superconducting phase dynamics provide additional control. Repeating these measurements across nominally identical devices would establish the reproducibility of the effect and separate device specific resonances from a general rectification mechanism.

        For the few-channel contacts, a microscopic Floquet calculation using the independently extracted pincodes and calibrated drive amplitudes would provide the most direct test of the threshold weighted PAMAR construction. Such a comparison could determine which remaining deviations originate from coherent interference between driven MAR paths, nonequilibrium occupations, or the electromagnetic environment. The phenomenological framework developed here remains useful for this purpose because it identifies the relevant charge orders and provides a transparent benchmark over complete amplitude maps.

        The subharmonic response at high conductance motivates a systematic study across contact transmission, microwave frequency, temperature, and magnetic field. Following the intermediate replicas as the contact evolves from the controlled few-channel regime into the high-conductance regime could establish where the additional \(p=2\) contribution emerges. Its amplitude dependence and response to suppressed superconductivity could distinguish correlated Cooper pair transfer from fractional phase locking or an environmentally generated replica process. These experiments would turn the present phenomenological charge assignment into a direct test of the microscopic origin of the subharmonic transport.
